% TOP_PROBLEM: {"id": 121, "title": "Sumsets of Perfect Squares", "category_id": 1, "category": {"id": 1, "name": "number_theory", "display_name": "Number Theory", "description": "Properties of integers, prime numbers, Diophantine equations.", "slug": "number-theory", "order_index": 1, "created_at": "2026-07-31T15:26:25.670Z"}, "status": "open"}
% FIRST_POSTED: null
% ENTRY_KIND: statement_only
% Imported from: unsolvedmath_additional_500.tex; UnsolvedMath ID 121
% !TeX program = lualatex
% Compile twice: lualatex 121.tex
% Self-contained: no external bibliography, images, or \input files.
% Numeric section labels and visible numbers are original UnsolvedMath IDs.
\documentclass[11pt,a4paper]{article}
\usepackage[margin=24mm,headheight=14pt]{geometry}
\usepackage{fontspec}
\setmainfont{Latin Modern Roman}
\setsansfont{Latin Modern Sans}
\setmonofont{Latin Modern Mono}
\usepackage{amsmath,amssymb,mathtools}
\usepackage{unicode-math}
\setmathfont{Latin Modern Math}
\usepackage{microtype}
\usepackage{enumitem,xurl,needspace}
\usepackage[dvipsnames]{xcolor}
\usepackage{hyperref}
\hypersetup{unicode=true,colorlinks=true,linkcolor=MidnightBlue,urlcolor=MidnightBlue,
 pdftitle={UnsolvedMath Problem 121},pdfauthor={Source-annotated compilation},bookmarksnumbered=true}
\usepackage{bookmark,tocloft,titlesec,fancyhdr}
\setlength{\cftsecnumwidth}{4.2em}
\cftsetpnumwidth{2.5em}
\cftsetrmarg{4em}
\titleformat{\section}{\Large\bfseries\raggedright}{\thesection}{0.7em}{}
\pagestyle{fancy}\fancyhf{}
\fancyhead[L]{\small\sffamily Additional UnsolvedMath statements}
\fancyhead[R]{\small Original catalogue IDs}
\fancyfoot[C]{\thepage}
\setlength{\parindent}{0pt}
\setlength{\parskip}{5pt plus 1pt}
\setlength{\emergencystretch}{3em}
\setcounter{tocdepth}{1}\setcounter{secnumdepth}{1}
\setlist[itemize]{leftmargin=3em,itemsep=4pt,topsep=4pt,parsep=0pt}
\allowdisplaybreaks
\newcommand{\N}{\mathbb{N}}
\newcommand{\Z}{\mathbb{Z}}
\newcommand{\Q}{\mathbb{Q}}
\newcommand{\R}{\mathbb{R}}
\newcommand{\C}{\mathbb{C}}
\newcommand{\F}{\mathbb{F}}
\newcommand{\A}{\mathbb{A}}
\newcommand{\PP}{\mathbb{P}}
\newcommand{\E}{\mathbb{E}}
\newcommand{\eps}{\varepsilon}
\newcommand{\dd}{\,\mathrm d}
\newcommand{\sref}[2]{\hyperlink{source:#1:#2}{[S#2]}}
\newif\ifshowcatalogue
\showcataloguetrue
\newcommand{\cataloguescope}[1]{\ifshowcatalogue
 \paragraph{Statement recorded in the supplied source.}
 #1\fi}
\begin{document}
% UnsolvedMath ID 121; source catalogue code GREEN-033

\setcounter{section}{120}

\section{Additive Growth of Sets of Squares}\label{121}

{\small\sffamily UnsolvedMath ID 121 · Number Theory}

\subsection{Definitions and mathematical statement}

\paragraph{Shared notation.}Unless a section says otherwise, $\N=\{1,2,\ldots\}$,
$\Z$ is the ring of integers, $\Q,\R,\C$ are the rational, real, and complex
fields, $[N]=\{1,\ldots,\lfloor N\rfloor\}$, and $\log$ is the natural
logarithm. For real functions of a parameter tending to infinity,
$f\ll g$ and $f=O(g)$ mean $|f|\leq Cg$ eventually for some fixed $C$;
$f\gg g$ reverses this comparison; $f=o(g)$ means $f/g\to0$;
$f\sim g$ means $f/g\to1$; and $f\asymp g$ means both order bounds.
A subscript on $O$ or $\ll$ specifies allowed constant dependence.
The expression $n^{o(1)}$ in an upper bound means growth smaller than
$n^\epsilon$ for each fixed $\epsilon>0$. Local definitions govern any
exception, finite-field convention, or use of the same symbol for another
object. An asymptotic claim is not automatically an effective algorithm.



A square here is a positive integer $u^2$, with $u\in\N$. For a finite set $A$ of such squares, $A+A=\{a+b:a,b\in A\}$ counts distinct sum values, not representations. The conjectured exponent $c>0$ must work uniformly over all finite $A$ with $|A|\geq2$, independently of the magnitudes of its elements. \sref{121}{1} \sref{121}{2}

\paragraph{Statement recorded in the supplied source.}
Is there an absolute constant $c > 0$ such that if $A \subset \mathbb{N}$ is a set of squares of size at least 2, then $|A + A| \geq |A|^{1+c}$? \sref{121}{1}

\paragraph{Residual task reported by the archive.}

The embedded review (2026-08-17) records: Prove a uniform fixed power gain, ideally identify the optimal c, or construct square sets with near-linear doubling that refute it. \sref{121}{1}

\subsection{Short English statement}

Do finite sets of perfect squares always gain a fixed power in size when all pairwise sums are formed?

\subsection{Sources}

{\small

\begin{itemize}

\item[\hypertarget{source:121:1}{S1}] ULAM.AI, \emph{UnsolvedMath}, supplied \texttt{problems.json}, record 121 (GREEN-033), statement and background. The embedded literature-review date is 2026-08-17; that is the archive’s review date, not a fresh certification. Dataset landing page: \url{https://huggingface.co/datasets/ulamai/UnsolvedMath}.

\item[\hypertarget{source:121:2}{S2}] M.-C. Chang, On problems of Erdős and Rudin, J. Functional Analysis 207 (2004), 444-460, \nolinkurl{doi:10.1016/S0022-1236}(03)00073-9. \url{https://doi.org/10.1016/S0022-1236(03)00073-9}. Bibliographic information imported from the supplied record.

\item[\hypertarget{source:121:3}{S3}] N. Hegyvári, Note on the sumset of squares, arXiv:2504.13230v2 (2025), withdrawal announced. \url{https://arxiv.org/abs/2504.13230}. Bibliographic information imported from the supplied record.

\item[\hypertarget{source:121:4}{S4}] B. Green, 100 Open Problems, Problem 60, version. \url{https://people.maths.ox.ac.uk/greenbj/papers/open-problems.pdf}. The author-hosted document was inspected on 27 September 2026; the archive’s local code is not a problem-number locator in this paper.

\end{itemize}

}

\label{end:121}
\end{document}
